% Save as: module9-problem-set.tex
% ============================================================================
% Problem Set 9: Digital & Platform Markets
% Microeconomic Theory for Experimentalists & Applied Microeconomists
% Built from templates/problem_set_template.tex. Compile with pdflatex.
% ============================================================================
\documentclass[11pt]{article}
\usepackage{amsmath, amssymb, amsthm}
\usepackage[margin=1in]{geometry}
\usepackage{enumitem}
\usepackage{booktabs}

\newtheorem{question}{Question}

\title{Problem Set 9: Digital \& Platform Markets\\
\large Microeconomic Theory for Experimentalists \& Applied Microeconomists}
\author{Due: \textbf{[date --- before Lecture 1 of Module 10]}}
\date{}

\begin{document}
\maketitle

\noindent\emph{Ground rules: collaboration on ideas is encouraged; write-ups
are individual. You may use AI tools for parts flagged with
$\langle$AI$\rangle$ and must attach the transcript. Show all calculations.}

\medskip

\noindent\textbf{Weights:} Part A --- 30 points. Part B --- 45 points.
Part C --- 25 points.

\bigskip

% ---------------------------------------------------------------------------
\section*{Part A: Theory --- Guided Calculations (30 points)}

\begin{question}[Two-sided pricing with numbers --- 15 points]
A rideshare platform faces rider demand $R(p) = 100 - 20p$ (riders per
day at rider fee $p$) and incurs a cost of \$3 per rider served.

\begin{enumerate}[label=(\alph*)]
  \item \textbf{(4 pts)} Treating the rider side in isolation, find the
        profit-maximizing fee and the resulting profit. \emph{Hint: the
        profit function is a parabola; its peak is midway between its
        zeros.}
  \item \textbf{(5 pts)} Now recognize the platform's other side: each
        active rider generates \$3 of profit on the driver side (higher
        utilization, commissions). The \$3 service cost still applies;
        the driver-side \$3 is additional profit per served rider. Find
        the new profit-maximizing rider fee and total platform profit
        (rider-side margin plus driver-side profit). State plainly what is remarkable about the
        optimal fee relative to the \$3 cost.
  \item \textbf{(3 pts)} A competition authority calls the fee in (b)
        predatory pricing. Write the platform's two-sentence defense
        using your arithmetic, and one sentence on what evidence would
        distinguish genuine predation from two-sided pricing.
  \item \textbf{(3 pts)} In one paragraph: which side of a two-sided
        platform should be the ``money side''? Use the
        elasticity-and-externality rule of thumb and one real example
        (cards, media, delivery apps).
\end{enumerate}
\end{question}

\begin{question}[Interference in a platform A/B test --- 15 points]
A rideshare market has fixed driver capacity of 1{,}000 completed rides
per day. Baseline: 1{,}000 ride requests per day, so all requests are
served; riders are split into two equal halves for an experiment
(500 baseline requests each). The treated half receives a discount that
raises its requests by 20\%; the control half's requests are unchanged.
When total requests exceed capacity, completed rides are allocated
proportionally to requests.

\begin{enumerate}[label=(\alph*)]
  \item \textbf{(5 pts)} Compute each group's completed rides under the
        experiment and the naive treatment estimate (treated minus
        control completed rides, and in percent of the 500-ride
        baseline).
  \item \textbf{(4 pts)} Now compute the all-treated counterfactual:
        every rider gets the discount. What happens to completed rides
        relative to baseline, and what is the true effect of full
        rollout on completed rides?
  \item \textbf{(3 pts)} Name the assumption that failed, say exactly
        where it failed (which resource do the groups share), and state
        the direction of the naive bias in one sentence.
  \item \textbf{(3 pts)} Two design fixes: cluster randomization and
        switchback randomization. One sentence each: what the fix buys,
        and the price it pays.
\end{enumerate}
\end{question}

% ---------------------------------------------------------------------------
\section*{Part B: Experimental Design (45 points)}

\begin{question}[Testing a marketplace feature under interference]
A large e-commerce marketplace wants to test a ``verified seller''
badge: sellers who pass a quality audit get a visible badge on their
listings. The concern: badges may not grow the market --- they may
merely reallocate demand from unbadged to badged sellers. Design the
experiment that separates market growth from reallocation. Specify:
\begin{enumerate}[label=(\alph*)]
  \item \textbf{(6 pts)} The setting and stakes: the marketplace
        structure (categories, sellers per category, buyer search
        behavior), why within-category demand stealing is the expected
        interference channel, and why the growth-vs-reallocation
        question determines whether the feature should ship.
  \item \textbf{(12 pts)} Treatments and randomization: your design
        must include a randomization unit chosen FOR the interference
        channel (justify cluster vs.\ switchback vs.\ a mixed design),
        and must identify (1) the total effect on badged sellers,
        (2) the spillover onto unbadged sellers in the same market,
        and (3) the market-level net effect. State precisely which
        comparison delivers each, and what a naive seller-level
        randomization would have concluded.
  \item \textbf{(9 pts)} Measurement: primary outcomes at seller and
        market level, guardrail metrics (Kohavi-style), novelty-effect
        handling (how long you run and why), and the
        sample-ratio-mismatch check.
  \item \textbf{(9 pts)} Hypotheses and power: state the
        growth-vs-reallocation hypotheses quantitatively (e.g.\ badged
        sellers +15\%, unbadged $-x$\%, market net $+y$\%), and give
        the cluster-level power reasoning: number of category-markets
        per arm given a stated intra-cluster correlation, and which
        comparison (the spillover or the net effect) drives the sample
        size and why.
  \item \textbf{(9 pts)} The biggest threat to YOUR design and your
        answer to it. (Candidates: buyers search across categories, so
        clusters leak; badge supply is endogenous --- sellers select
        into auditing; switchback carryover through reviews and repeat
        purchase.)
\end{enumerate}
\end{question}

% ---------------------------------------------------------------------------
\section*{Part C: Silicon Pilot Interpretation $\langle$AI$\rangle$
(25 points)}

\begin{question}[The sellers that colluded in three rounds]
You run \texttt{module9-simulation-lab.ipynb}. (For reference: the
lab's prompts hand each agent a plain-language description of demand
and costs; \texttt{history} agents see the full price/profit path each
round, \texttt{amnesia} agents are stateless; and the lab records only
prices --- no reasoning traces.) In the history
treatment, your LLM seller pairs converge to the joint-monopoly price
of 6 within three rounds and hold it there, with no communication.
Calvano et al.'s (2020) Q-learning agents needed on the order of
100{,}000 learning episodes to reach supra-competitive prices. Your
lab partner drafts an abstract: ``LLM pricing agents collude faster
than reinforcement learners --- deployed LLM pricing tools pose an
immediate collusion risk.''

\begin{enumerate}[label=(\alph*)]
  \item \textbf{(8 pts)} Explain why three-round convergence to
        exactly the monopoly price is NOT evidence of ``learning to
        collude'' in Calvano's sense, and give two mechanisms (at
        least one LLM-specific) that could produce it. \emph{Hint:
        what did your agents receive in the prompt that Calvano's
        Q-learners had to discover through experience --- and what has
        your model read?}
  \item \textbf{(9 pts)} Design the diagnostic battery: your design
        must include (i) the \emph{hidden-demand} run (agents observe
        only their own realized sales, never the demand function ---
        the information condition of deployed pricing tools) and
        (ii) the \emph{manual-deviation probe} (force one agent to cut
        price for one round; Calvano's signature is
        punishment-then-forgiveness). Add at least one more (shifted
        cost/benchmark numbers; three-seller markets). State the
        predicted result pattern under each mechanism from (a). Run it
        if API budget allows and report; otherwise state predictions
        precisely.
  \item \textbf{(8 pts)} The twist (Module 6's lesson returns): argue
        that a properly designed version of this simulation IS a
        legitimate, policy-relevant study --- of what, exactly? State
        the research question it answers, the one design condition
        that must hold for policy relevance (hint: information), and
        the one-sentence scope limitation for the abstract your
        partner should write instead.
\end{enumerate}
Attach your AI transcripts and, if you ran code, the results CSV.
\end{question}

\end{document}
